\documentclass{article}
\usepackage[T1]{fontenc}
\usepackage{lmodern}
\usepackage{graphicx}
\usepackage{amsmath,amssymb}
\usepackage[a4paper,margin=2cm]{geometry}
\usepackage[absolute,overlay]{textpos}
\setlength{\TPHorizModule}{1mm}
\setlength{\TPVertModule}{1mm}
\usepackage[indonesian]{babel}
\usepackage{hyperref}
\setlength{\emergencystretch}{2em}
% unit: Halaman 58

\begin{document}

% === Halaman 58 (llm) ===

(ii) \begin{align*}
a \cdot b &= (x^6 + x^4 + x^2 + x + 1) \cdot (x^7 + x + 1) \\
&= x^{13} + x^7 + x^6 + x^{11} + x^5 + x^4 + x^9 + x^3 + x^2 + x^8 + x^2 + x + x^7 + x + 1 \pmod{2} \\
&= x^{13} + x^{11} + x^9 + x^8 + 2x^7 + x^6 + x^5 + x^4 + x^3 + 2x^2 + 2x + 1 \pmod{2} \\
&= x^{13} + x^{11} + x^9 + x^8 + 0x^7 + x^6 + x^5 + x^4 + x^3 + 0x^2 + 0x + 1 \\
&= x^{13} + x^{11} + x^9 + x^8 + x^6 + x^5 + x^4 + x^3 + 1
\end{align*}

Karena $m = 8$, maka hasil perkalian di atas direduksi menjadi derajat $< 8$ oleh \textit{irreducible polynomial}. Algoritma Rijndael menggunakan \textit{irreducible polynomial} $f(x) = x^8 + x^4 + x^3 + x + 1$

\[ x^{13} + x^{11} + x^9 + x^8 + x^6 + x^5 + x^4 + x^3 + 1 \pmod{f(x)} = x^7 + x^6 + 1 = 11000001 \]

$\dots a \cdot b = 11000001 = \text{'C1'}$

\end{document}
